\begin{flushleft}
    \footnotesize \textbf{ÁREAS DE ESTUDIO:} \textit{agentes de IA, hiperautomatización, WhatsApp Business, n8n, arquitectura de software, APIs REST, gobierno de IA, trazabilidad, PyMEs, automatización de pedidos.}
\end{flushleft}

\begin{abstract}
DeliverAI es un proyecto académico aplicado que busca reducir la fricción entre una consulta comercial por WhatsApp y un pedido confirmado, completo y trazable. El caso se enfoca en PyMEs que reciben pedidos personalizados por chat, donde la información llega en texto libre, imágenes, cambios de último minuto y datos incompletos. La solución propuesta combina un agente conversacional orquestado por n8n, una API transaccional en Spring Boot, una base PostgreSQL y una interfaz administrativa en React. El artículo presenta el problema, el alcance real del MVP técnico, la arquitectura defendible, el valor esperado, los riesgos críticos y los controles necesarios para avanzar a un piloto controlado. La tesis central es prudente: DeliverAI no debe presentarse como producción terminada, sino como un prototipo funcional con evidencia suficiente para validar, con métricas y supervisión humana, si un agente de IA puede mejorar la conversión de conversaciones de WhatsApp a pedidos confirmados.
\end{abstract}

\epigraph{\textit{``Un agente no está listo para producción cuando funciona; está listo cuando puede limitarse, monitorearse, explicarse, auditarse, detenerse y corregirse.''}}{--- Principio de gobierno usado en el curso}
